\documentclass[a4paper,12pt]{article}

% ======== Packages ========
\usepackage[top=1in,bottom=1in,left=1in,right=1in]{geometry}

% ---------- 断行救急：防止行内不可断 token（代码路径/编号/±20% 串）硬越界 ----------
% ctex 中文两端对齐的字符间胶水是刚性的，行尾挤不下不可断 token 时 LaTeX 会硬凸出行宽；
% \emergencystretch 给断行器 2em 额外弹性，在判定 overfull 前先收紧相邻行。零副作用。
% 对超出 2em 的巨型 token（如 P0393/P0526 长编号串），写作时手动插断点：
%   $\{P0393,\allowbreak P0526\}$、\texttt{code/}\allowbreak 目录
\emergencystretch=2em

\setlength{\parindent}{2em}
\usepackage{amsmath,amssymb}
\usepackage{fontspec}
\usepackage{titlesec}
\usepackage{fancyhdr}
\usepackage{booktabs}
\usepackage{listings}
\usepackage{graphicx}
\usepackage{ifthen}
\usepackage{times}
\usepackage{hyperref}

% ======== Fonts ========
\setmainfont{Times New Roman}
\setsansfont{Arial}
\setmonofont{Courier New}

% ======== Variables ========
\newcommand{\MCMYear}{2026}
\newcommand{\TeamNumber}{0000}
\newcommand{\ProblemChosen}{A}
\newcommand{\PaperTitle}{[Paper Title]}
\newcommand{\PaperDate}{June 3, 2026}
\newcommand{\SummaryText}{[Summary: problem overview + methods for each sub-problem + key numerical results + conclusions, 300-500 words]}
\newcommand{\KeywordsText}{[Keyword1]; [Keyword2]; [Keyword3]}
\newcommand{\TotalPages}{4}
\newcommand{\kw}[1]{\textbf{#1}} % content emphasis: conclusion phrase > model/method name > key number in phrase

% ======== Page Header ========
\fancyhf{}
\fancyhead[L]{\sffamily Team \#\TeamNumber}
\fancyhead[R]{\sffamily Page \thepage\ of \TotalPages}
\fancyhead[C]{\rule[-0.1em]{0pt}{\baselineskip}}
\cfoot{}
\pagestyle{fancy}

% ======== Heading Formats ========
\titleformat{\section}[block]{\normalsize\bfseries}{\thesection\ }{0em}{}
\titleformat{\subsection}[block]{\normalsize\bfseries}{\thesubsection\ }{0em}{}
\titlespacing*{\section}{0pt}{1.2em}{0.55em}
\titlespacing*{\subsection}{0pt}{0.82em}{0.4em}

% ======== Three-line Table ========
\newcommand{\threelinetable}[4]{%
  \begin{table}[htbp]
    \centering
    \caption{#1}
    \begin{tabular}{#2}
      \toprule
      #3 \\
      \midrule
      #4 \\
      \bottomrule
    \end{tabular}
  \end{table}%
}

% ======== Code Block ========
\lstset{
  frame=topbottom,
  framerule=0.55pt,
  framesep=0pt,
  basicstyle=\ttfamily\small,
  backgroundcolor=\color[gray]{0.97},
  numbers=none,
  showstringspaces=false,
  breaklines=true,
}

% ======== Table of Contents Entry ========
\newcommand{\tocentry}[4]{\par\vskip0.55em\noindent%
  \ifthenelse{\equal{#4}{1}}{\bfseries}{}%
  #1\quad#2%
  \leaders\hbox{$\cdot$}\hfill\ #
  #3\par%
}

% ======== Document ========
\begin{document}

% ---- Summary Sheet ----
\thispagestyle{empty}
\vspace*{-1.8em}
\begin{center}
  \begin{tabular}{ccc}
    \textbf{Problem Chosen} & \textbf{\MCMYear MCM/ICM Summary Sheet} & \textbf{Team Control Number} \\[0.7em]
    \textbf{\LARGE\ProblemChosen} & & \textbf{\LARGE\TeamNumber}
  \end{tabular}
\end{center}
\vspace{0.43em}
\hrule height 0.8pt
\vspace{2.4em}
\begin{center}\textbf{Summary}\end{center}
\vspace{2.1em}
\SummaryText
\vspace{0.75em}
\par\textbf{Keywords:} \KeywordsText
\newpage

% ---- Main Content ----
\pagestyle{fancy}
\input{sections/1_introduction}
\pagebreak[100]
\input{sections/2_assumptions}
\pagebreak[100]
\input{sections/3_model_design}
\pagebreak[100]
\input{sections/4_solution}
\input{sections/5_sensitivity}
\input{sections/6_strengths_weaknesses}
\pagebreak[100]
\input{sections/7_conclusions}

% ---- References ----
\pagebreak[100]
\fontsize{20pt}{24pt}\selectfont\textbf{Use of AI}\par\normalsize
\textbf{We used AI tools as follows:} \textit{[state tool name/version, purpose and stages, prompt approach, and how outputs were reviewed/verified; delete if unused]}%
\par\vspace{1em}
\fontsize{20pt}{24pt}\selectfont\textbf{References}\par\normalsize
\input{references}

% ---- Appendices ----
\pagebreak[100]
\fontsize{28pt}{32pt}\selectfont\textbf{Appendices}\par
\fontsize{20pt}{24pt}\selectfont\textbf{Appendix A \quad Core Code}\par\normalsize
\input{sections/A_code}

\end{document}
